\section{AIME II 2010}
        \begin{enumerate}
        	\item (\textit{AIME II 2010}) Let $N$ be the greatest integer multiple of $36$ all of whose digits are even and no two of whose digits are the same. Find the remainder when $N$ is divided by $1000$.


 
 \textbf{Answer: }640



\textbf{Solution}

If an integer is divisible by $36$, it must also be divisible by $9$ since $9$ is a factor of $36$. It is a well-known fact that, if $N$ is divisible by $9$, the sum of the digits of $N$ is a multiple of $9$. Hence, if $N$ contains all the even digits, the sum of the digits would be $0 + 2 + 4 + 6 + 8 = 20$, which is not divisible by $9$ and thus $36$.The next logical try would be $8640$, which happens to be divisible by $36$. Thus, $N = 8640 \equiv \boxed{640} \pmod {1000}$.

 


 




	\item (\textit{AIME II 2010}) A point $P$ is chosen at random in the interior of a unit square $S$. Let $d(P)$ denote the distance from $P$ to the closest side of $S$. The probability that $\frac{1}{5}\le d(P)\le\frac{1}{3}$ is equal to $\frac{m}{n}$, where $m$ and $n$ are relatively prime positive integers. Find $m+n$.


 
 \textbf{Answer: }281



\textbf{Solution}

Any point outside the square with side length $\frac{1}{3}$ that has the same center and orientation as the unit square and inside the square with side length $\frac{3}{5}$ that has the same center and orientation as the unit square has $\frac{1}{5}\le d(P)\le\frac{1}{3}$.

 
                    \begin{figure}[H]
                        \centering
                        \includegraphics[width=0.5\textwidth]{Solution/AIME_II_2010_2_Solution_Figure_1.png}
                    \end{figure}
                    Since the area of the unit square is $1$, the probability of a point $P$ with $\frac{1}{5}\le d(P)\le\frac{1}{3}$ is the area of the shaded region, which is the difference of the area of two squares.

 $\left(\frac{3}{5}\right)^2-\left(\frac{1}{3}\right)^2=\frac{9}{25}-\frac{1}{9}=\frac{56}{225}$

 Thus, the answer is $56 + 225 = \boxed{281}.$

 





\textbf{Solution 2}

First, let's figure out $d(P) \geq \frac{1}{3}$ which is
 \[\left(\frac{3}{5}\right)^2=\frac{9}{25}.\]Then, $d(P) \geq \frac{1}{5}$ is a square inside $d(P) \geq \frac{1}{3}$, so
 \[\left(\frac{1}{3}\right)^2=\frac{1}{9}.\]Therefore, the probability that $\frac{1}{5}\le d(P)\le\frac{1}{3}$ is
 \[\frac{9}{25}-\frac{1}{9}=\frac{56}{225}\]So, the answer is $56+225=\boxed{281}$

 





\textbf{Solution 3}

First, lets assume that point $P$ is closest to a side $S$ of the square. If it is $\frac{1}{5}$ far from $S$, then it should be at least $\frac{1}{5}$ from both the adjacent sides of $S$ in the square. This leaves a segment of $1 - 2 \cdot \frac{1}{5} = \frac{3}{5}$. If the distance from $P$ to $S$ is $\frac{1}{3}$, then notice the length of the side-ways segment for $P$ is $1 - 2 \cdot \frac{1}{3} = \frac{1}{3}$. Notice that as the distance from $P$ to $S$ increases, the possible points for the side-ways decreases. This produces a trapezoid with parallel sides $\frac{3}{5}$ and $\frac{1}{3}$ with height $\frac{1}{3} - \frac{1}{5} = \frac{2}{15}$. This trapezoid has area (or probability for one side) $\frac{1}{2} \cdot \left(\frac{1}{3}+\frac{3}{5}\right)\cdot \frac{2}{15} = \frac{14}{225}$. Since the square has $4$ sides, we multiply by $4$. Hence, the probability is $\frac{56}{225}$. The answer is $\boxed{281}$.               ~Saucepan\_man02

 




	\item (\textit{AIME II 2010}) Let $K$ be the product of all factors $(b-a)$ (not necessarily distinct) where $a$ and $b$ are integers satisfying $1\le a < b \le 20$. Find the greatest positive integer $n$ such that $2^n$ divides $K$.


 
 \textbf{Answer: }150



\textbf{Solution}

In general, there are $20-n$ pairs of integers $(a, b)$ that differ by $n$ because we can let $b$ be any integer from $n+1$ to $20$ and set $a$ equal to $b-n$. Thus, the product is $(1^{19})(2^{18})\cdots(19^1)$ (or alternatively, $19! \cdot 18! \cdots 1!$.)

 When we count the number of factors of $2$, we have 4 groups, factors that are divisible by $2$ at least once, twice, three times and four times.

 


 
                \begin{itemize}
                    \item Numbers that are divisible by $2$ at least once: $2, 4, \cdots, 18$
                \end{itemize}


                
            \begin{enumerate}[label=(\roman*)]
                \item Exponent corresponding to each one of them $18, 16, \cdots 2$
            \end{enumerate}
            
            \begin{enumerate}[label=(\roman*)]
                \item Sum $=2+4+\cdots+18=\frac{(20)(9)}{2}=90$
            \end{enumerate}
            


 
                \begin{itemize}
                    \item Numbers that are divisible by $2$ at least twice: $4, 8, \cdots, 16$
                \end{itemize}


                
            \begin{enumerate}[label=(\roman*)]
                \item Exponent corresponding to each one of them $16, 12, \cdots 4$
            \end{enumerate}
            
            \begin{enumerate}[label=(\roman*)]
                \item Sum $=4+8+\cdots+16=\frac{(20)(4)}{2}=40$
            \end{enumerate}
            


 
                \begin{itemize}
                    \item Numbers that are divisible by $2$ at least three times: $8,16$
                \end{itemize}


                
            \begin{enumerate}[label=(\roman*)]
                \item Exponent corresponding to each one of them $12, 4$
            \end{enumerate}
            
            \begin{enumerate}[label=(\roman*)]
                \item Sum $=12+4=16$
            \end{enumerate}
            


 
                \begin{itemize}
                    \item Number that are divisible by $2$ at least four times: $16$
                \end{itemize}


                
            \begin{enumerate}[label=(\roman*)]
                \item Exponent corresponding to each one of them $4$
            \end{enumerate}
            
            \begin{enumerate}[label=(\roman*)]
                \item Sum $=4$
            \end{enumerate}
            


 Summing these give an answer of $\boxed{150}$.

 




	\item (\textit{AIME II 2010}) Dave arrives at an airport which has twelve gates arranged in a straight line with exactly $100$ feet between adjacent gates. His departure gate is assigned at random. After waiting at that gate, Dave is told the departure gate has been changed to a different gate, again at random. Let the probability that Dave walks $400$ feet or less to the new gate be a fraction $\frac{m}{n}$, where $m$ and $n$ are relatively prime positive integers. Find $m+n$.


 
 \textbf{Answer: }052



\textbf{Solution 1}

There are $12 \cdot 11 = 132$ possible situations ($12$ choices for the initially assigned gate, and $11$ choices for which gate Dave's flight was changed to). We are to count the situations in which the two gates are at most $400$ feet apart.

 If we number the gates $1$ through $12$, then gates $1$ and $12$ have four other gates within $400$ feet, gates $2$ and $11$ have five, gates $3$ and $10$ have six, gates $4$ and $9$ have have seven, and gates $5$, $6$, $7$, $8$ have eight. Therefore, the number of valid gate assignments is 
 \[2\cdot(4+5+6+7)+4\cdot8 = 2 \cdot 22 + 4 \cdot 8 = 76\] so the probability is $\frac{76}{132} = \frac{19}{33}$. The answer is $19 + 33 = \boxed{052}$.

 





\textbf{Solution 2}

As before, derive that there are $132$ possibilities for Dave's original and replacement gates.

 Now suppose that Dave has to walk $100k$ feet to get to his new gate.  This means that Dave's old and new gates must be $k$ gates apart.  (For example, a $100$ foot walk would consist of the two gates being adjacent to each other.)  There are $12-k$ ways to pick two gates which are $k$ gates apart, and $2$ possibilities for gate assignments, for a total of $2(12-k)$ possible assignments for each $k$.

 As a result, the total number of valid gate arrangements is 
 \[2\cdot 11 + 2\cdot 10 + 2\cdot 9 + 2\cdot 8 = 76\] and so the requested probability is $\frac{19}{33}$ for a final answer of $\boxed{052}$.

 





\textbf{Solution 3}

We can just manually count via each probability.

 Note that the probability he gets assigned to gate \( 1 \) is the same as the probability he gets assigned to gate \( 12 \). Mappings are made symmetrically. Thus we find the cases until gate \( 6 \), and multiply each by \( 2 \).

 If he gets gate \( 1 \) with chance \( \frac{1}{12} \), he can move to gate \( 2, 3, 4, \) or \( 5 \) with probability \( \frac{4}{12} \) after the change. The total probability is \( \frac{1}{12} \times \frac{4}{11} \).

 If he gets gate \( 2 \) with chance \( \frac{1}{12} \), he can move to gate \( 1, 3, 4, 5, \) or \( 6 \), with combined probability \( \frac{5}{12} \). The total probability is \( \frac{1}{12} \times \frac{5}{11} \).

 If he gets gate \( 3 \) with chance \( \frac{1}{12} \), he can move to \( 1, 2, 4, 5, 6, 7 \) with combined probability \( \frac{6}{12} \). The total probability is \( \frac{1}{12} \times \frac{6}{11} \).

 If he gets gate \( 4 \) with chance \( \frac{1}{12} \), he can move to \( 1, 2, 3, 5, 6, 7, 8 \) with combined probability \( \frac{7}{12} \). The total probability is \( \frac{1}{12} \times \frac{7}{11} \).

 If he gets gate \( 5 \) with chance \( \frac{1}{12} \), he can move to \( 1, 2, 3, 4, 6, 7, 8, 9 \) with combined probability \( \frac{8}{12} \). The total probability is \( \frac{1}{12} \times \frac{8}{11} \).

 Finally, if he gets gate \( 6 \) with chance \( \frac{1}{12} \), he can move to \( 2, 3, 4, 5, 7, 8, 9, 10 \) with combined probability \( \frac{8}{12} \). The total probability is \( \frac{1}{12} \times \frac{8}{11} \).

 Then, the combined probability for \( 1 \text{--} 6 \) is:
 \[\frac{1}{12} \times \frac{4}{11} + \frac{1}{12} \times \frac{5}{121} + \frac{1}{12} \times \frac{6}{11} + \frac{1}{12} \times \frac{7}{11} + \frac{1}{12} \times \frac{8}{11} + \frac{1}{12} \times \frac{8}{11}\]or \( \frac{1}{12} \left( \frac{4}{11} + \frac{5}{11} + \frac{6}{11} + \frac{7}{11} + \frac{8}{11} +  \frac{8}{11} \right) \), or \( \frac{1}{132} (4 + 5 + 6 + 7 + 8 + 8) \), or \( \frac{1}{132} (38) \) or \( \frac{38}{132} \), or \( \frac{19}{66} \).

 Then, when we multiply the total by \( 2 \), we obtain \( \frac{19}{33} \). $19+33 = \boxed{052}$

 





\textbf{Solution 4 (fast)}

We can count the favorable outcomes. Note that it doesn't matter which is the starting and ending gate, only the gap between them matters. 

 For an outcome to be favorable, the two gates must have a gap \( \leq 4 \). There are 11 ways to have gap 1, 10 for gap 2, 9 for gap 3, and 8 for gap 4, giving us \( 11 + 10 + 9 + 8 = 38 \) favorable outcomes. There are \( \binom{12}{2} = 66 \) total outcomes, and so the probability of getting a favorable outcome is \( \frac{38}{66} = \frac{19}{33} \). $19+33=\boxed{052}$.

 




	\item (\textit{AIME II 2010}) Positive numbers $x$, $y$, and $z$ satisfy $xyz = 10^{81}$ and $(\log_{10}x)(\log_{10} yz) + (\log_{10}y) (\log_{10}z) = 468$. Find $\sqrt {(\log_{10}x)^2 + (\log_{10}y)^2 + (\log_{10}z)^2}$.


 
 \textbf{Answer: }075



\textbf{Solution}

Using the properties of logarithms, $\log_{10}xyz = 81$ by taking the log base 10 of both sides, and $(\log_{10}x)(\log_{10}y) + (\log_{10}x)(\log_{10}z) + (\log_{10}y)(\log_{10}z)= 468$ by using the fact that $\log_{10}ab = \log_{10}a + \log_{10}b$. 

 Through further simplification, we find that $\log_{10}x+\log_{10}y+\log_{10}z = 81$. It can be seen that there is enough information to use the formula $\ (a+b+c)^{2} = a^{2}+b^{2}+c^{2}+2ab+2ac+2bc$, as we have both $\ a+b+c$ and $\ 2ab+2ac+2bc$, and we want to find $\sqrt {a^2 + b^2 + c^2}$. 

 After plugging in the values into the equation, we find that $\ (\log_{10}x)^2 + (\log_{10}y)^2 + (\log_{10}z)^2$ is equal to $\ 6561 - 936 = 5625$. 

 However, we want to find $\sqrt {(\log_{10}x)^2 + (\log_{10}y)^2 + (\log_{10}z)^2}$, so we take the square root of $\ 5625$, or $\boxed{075}$.

 





\textbf{Solution 2}

Let $a=\log_{10}x$, $b=\log_{10}y,$ and $c=\log_{10}z$.

 We have $a+b+c=81$ and $a(b+c)+bc=ab+ac+bc=468$. Since these two equations look a lot like Vieta's for a cubic, create the polynomial $x^3-81x^2+468x=0$ (leave the constant term as $0$ to make things easy). Dividing by $x$ yields $x^2-81x+468=0$.

 Now we use the quadratic formula:

 $x=\frac{81\pm\sqrt{81^2-4\cdot468}}{2}$

 $x=\frac{81\pm\sqrt{4689}}{2}$

 $x=\frac{81+3\sqrt{521}}{2}$, $x=\frac{81-3\sqrt{521}}{2}$

 Since the question asks for $\sqrt{a^2+b^2+c^2}$  (remember one of the values was the solution $x=0$ that we divided out in the beginning), we find:

 $\sqrt{\left(\frac{81+3\sqrt{521}}{2}\right)^2+\left(\frac{81-3\sqrt{521}}{2}\right)^2}$

 $=\sqrt{2\cdot\frac{81^2+(3\sqrt{521})^2}{4}}$

 $=\sqrt{\frac{11250}{2}}$

 $=\boxed{075}$

 




	\item (\textit{AIME II 2010}) Find the smallest positive integer $n$ with the property that the polynomial $x^4 - nx + 63$ can be written as a product of two nonconstant polynomials with integer coefficients.


 
 \textbf{Answer: }008



\textbf{Solution 1}

You can factor the polynomial into two quadratic factors or a linear and a cubic factor.

 For two quadratic factors, let $x^2+ax+b$ and $x^2+cx+d$ be the two quadratics, so that

 
 \[(x^2 + ax + b )(x^2 + cx + d) = x^4 + (a + c)x^3 + (b + d + ac)x^2 + (ad + bc)x + bd.\]

 Therefore, again setting coefficients equal, $a + c = 0\Longrightarrow a=-c$, $b + d + ac = 0\Longrightarrow b+d=a^2$ , $ad + bc = - n$, and so $bd = 63$.

 Since $b+d=a^2$, the only possible values for $(b,d)$ are $(1,63)$ and $(7,9)$. From this we find that the possible values for $n$ are $\pm 8 \cdot 62$ and $\pm 4 \cdot 2$. 

 For the case of one linear and one cubic factor, doing a similar expansion and matching of the coefficients gives the smallest $n$ in that case to be $48$. 

 Therefore, the answer is $4 \cdot 2 = \boxed{008}$.

 





\textbf{Solution 2}

Let $x^4-nx+63=(x^2+ax+b)(x^2+cx+d)$. From this, we get that $bd=63\implies d=\frac{63}{b}$ and $a+c=0\implies c=-a$. Plugging this back into the equation, we get $x^4-nx+63=(x^2+ax+b)\left(x^2-ax+\frac{63}{b}\right)$. Expanding gives us $x^4-nx+63=x^4+\left(-a^2+b+\frac{63}{b}\right)x^2+\left(\frac{63a}{b}-ab\right)x+63$. Therefore $-a^2+b+\frac{63}{b}=0$. Simplifying gets us $b(a^2-b)=63$. Since $a$ and $b$ must be integers, we can use guess and check for values of $b$ because $b$ must be a factor of $63$. Note that $b$ cannot be negative because $a$ would be imaginary. After guessing and checking, we find that the possible values of $(a,b)$ are $(\pm 8, 1), (\pm 4, 7), (\pm 4, 9),$ and $(\pm 8, 63)$. We have that $n=ab-\frac{63a}{b}$. Plugging in our values for $a$ and $b$, we get that the smallest positive integer value $n$ can be is $\boxed{008}$.-Heavytoothpaste

 




	\item (\textit{AIME II 2010}) Let $P(z)=z^3+az^2+bz+c$, where $a$, $b$, and $c$ are real. There exists a complex number $w$ such that the three roots of $P(z)$ are $w+3i$, $w+9i$, and $2w-4$, where $i^2=-1$. Find $|a+b+c|$.


 
 \textbf{Answer: }136



\textbf{Solution 1 (Vieta's)}

Set $w=x+yi$, so $x_1 = x+(y+3)i$, $x_2 = x+(y+9)i$, $x_3 = 2x-4+2yi$.

 Since $a,b,c\in{R}$, the imaginary part of $a,b,c$ must be $0$.

 Start with a, since it's the easiest one to do: $y+3+y+9+2y=0, y=-3$,

 and therefore: $x_1 = x$, $x_2 = x+6i$, $x_3 = 2x-4-6i$.

 Now, do the part where the imaginary part of c is 0 since it's the second easiest one to do: $x(x+6i)(2x-4-6i)$. The imaginary part is $6x^2-24x$, which is 0, and therefore $x=4$, since $x=0$ doesn't work.

 So now, $x_1 = 4, x_2 = 4+6i, x_3 = 4-6i$,

 and therefore: $a=-12, b=84, c=-208$. Finally, we have $|a+b+c|=|-12+84-208|=\boxed{136}$.

 





\textbf{Solution 1b}

Same as solution 1 except that when you get to $x_1 = x$, $x_2 = x+6i$, $x_3 = 2x-4-6i$, you don't need to find the imaginary part of $c$. We know that $x_1$ is a real number, which means that $x_2$ and $x_3$ are complex conjugates. Therefore, $x=2x-4$.

 





\textbf{Solution 2 (casework)}

Note that at least one of $w+3i$, $w+9i$, or $2w-4$ is real by complex conjugate roots. We now separate into casework based on which one.

 Let $w=x+yi$, where $x$ and $y$ are reals.

 Case 1: $w+3i$ is real. This implies that $x+yi+3i$ is real, so by setting the imaginary part equal to zero we get $y=-3$, so $w=x-3i$. Now note that since $w+3i$ is real, $w+9i$ and $2w-4$ are complex conjugates. Thus $\overline{w+9i}=2w-4$, so $\overline{x+6i}=2(x-3i)-4$, implying that $x=4$, so $w=4-3i$.

 Case 2: $w+9i$ is real. This means that $x+yi+9i$ is real, so again setting imaginary part to zero we get $y=-9$, so $w=x-9i$. Now by the same logic as above $w+3i$ and $2w-4$ are complex conjugates. Thus $\overline{w+3i}=2w-4$, so $\overline{x-6i}=2(x-9i)-4$, so $x+6i=2x-4-18i$, which has no solution as $x$ is real.

 Case 3: $2w-4$ is real. Going through the same steps, we get $y=0$, so $w=x$. Now $w+3i$ and $w+6i$ are complex conjugates, but $w=x$, which means that $\overline{x+3i}=x+6i$, so $x-3i=x+6i$, which has no solutions.

 Thus case 1 is the only one that works, so $w=4-3i$ and our polynomial is $(z-(4))(z-(4+6i))(z-(4-6i))$. Note that instead of expanding this, we can save time by realizing that the answer format is $|a+b+c|$, so we can plug in $z=1$ to our polynomial to get the sum of coefficients, which will give us $a+b+c+1$. Plugging in $z=1$ into our polynomial, we get $(-3)(-3-6i)(-3+6i)$ which evaluates to $-135$. Since this is $a+b+c+1$, we subtract 1 from this to get $a+b+c=-136$, so $|a+b+c|=\boxed{136}$.

 





\textbf{Solution 3}

By Vieta's we know the sum of the roots must be $-a$, a real number. That means $4w+12i-4$ is a real number, meaning $w$ has an imaginary component of $-3i$.

 Now we write $w = x-3i$. Then, $w+3i$ is the real root, meaning the other two are complex conjugates. We have $\overline{x+6i} = 2x-4-6i$, and solving, we get $x=4$. Then, $f(x) = (x-4)(x-4-6i)(x-4+6i) = (x-4)(x^2-8x+52)$. 

 We get $|a+b+c| = |-12+84-208| = \boxed{136}$.

 
-skibbysiggy

 




	\item (\textit{AIME II 2010}) Let $N$ be the number of ordered pairs of nonempty sets $\mathcal{A}$ and $\mathcal{B}$ that have the following properties:


 
                \begin{itemize}
                    
\item  $\mathcal{A} \cup \mathcal{B} = \{1,2,3,4,5,6,7,8,9,10,11,12\}$,
\item  $\mathcal{A} \cap \mathcal{B} = \emptyset$, 
\item  The number of elements of $\mathcal{A}$ is not an element of $\mathcal{A}$,
\item  The number of elements of $\mathcal{B}$ is not an element of $\mathcal{B}$.

                \end{itemize}


                Find $N$.


 
 \textbf{Answer: }772



\textbf{Solution}

Let us \href{https://artofproblemsolving.com/wiki/index.php?title=Partition}{partition} the set $\{1,2,\cdots,12\}$ into $n$ numbers in $A$ and $12-n$ numbers in $B$, 

 Since $n$ must be in $B$ and $12-n$ must be in $A$ ($n\ne6$, we cannot partition into two sets of 6 because $6$ needs to end up somewhere, $n\ne 0$ or $12$ either).

 We have $\dbinom{10}{n-1}$  ways of picking the numbers to be in $A$. 

 So the answer is $\left(\sum_{n=1}^{11} \dbinom{10}{n-1}\right) - \dbinom{10}{5}=2^{10}-252= \boxed{772}$.

 Note: We have $\dbinom{10}{n-1}$  ways of picking the numbers to be in $A$ because there are $n$ numbers in $A$ and since $12-n$ is already a term in the set we simply have to choose another $n-1$ numbers from the $10$ numbers that are available.

 





\textbf{Solution 2}

Regardless of the size $n$ of $A$ (ignoring the case when $n = 6$), $n$ must not be in $A$ and $12 - n$ must be in $A$. 

 There are $10$ remaining elements whose placements have yet to be determined. Note that the actual value of $n$ does not matter; there is always $1$ necessary element, $1$ forbidden element, and $10$ other elements that need to be distributed. There are $2$ places to put each of these elements, for $2^{10}$ possibilities.

 However, there is the edge case of $n = 6; 6$ is forced not to be in either set, so we must subtract the $\dbinom{10}{5}$ cases where $A$ and $B$ have size $6$.

 Thus, our answer is $2^{10} - \dbinom{10}{5} = 1024 - 252 = \boxed{772}$

 





\textbf{Solution 3 (PIE and Complementary Counting)}

The total number of possible subsets is $\sum_{i=1}^{11}\dbinom{12}{i}$, which is $2^{12}-2$. Note that picking a subset from the set leaves the rest of the set to be in the other subset. We exclude $i=0$ and $i=12$ since they leave a set empty. We proceed with complementary counting and casework:

 We apply the Principle of Inclusion and Exclusion for casework on the complementary cases. We find the ways where $|A|$ is in $A$, which violates the first condition. Then we find the ways where the elements $|B|$ and $12-|B|$ are in set $B$, which violates only the second condition, and not the first. This ensures we do not overcount.

 Case 1: $|A|$ is an element in $A$

 There are $\sum_{i=1}^{11}\dbinom{11}{i-1}$ = $2^{11}-1$ ways in this case.

 Case 2: $|B|$ and $12-|B|$ are in $B$

 We introduce a subcase where $|B|$ is not 6, since the other element would also be 6.There are $B-2$ elements to choose from 10 elements. Therefore, $|B|$ can be from 2 to 11.There are $\sum_{i=2}^{11}\dbinom{10}{i-2}-\dbinom{10}{4} = 2^{10}-211$ ways.The other subcase is when $|B|$ is equal to 6. There are $\dbinom{11}{5}=462$ ways.Adding the subcases gives us $2^{10}+251$. 

 Adding this with case one gives us $2^{11}+2^{10}+250$. Subtracting this from $2^{12}-2$ gives $1024-252=\boxed{772}$.

 





\textbf{Solution 4 (Not the best solution- Casework)}

Case 1: Set A has 1 element. Set B has 11 elements. 11 must be in set A, so there is $1$ way.

 Case 2: Set A has 2 elements. Set B has 10 elements. 10 must be in set A, and 2 must not be in set A, so there is $1\cdot\binom{10}{1}$ (for the remaining element).

 Case 3: Set A has 3 elements. Set B has 9 elements. 9 must be in set A, and 3 must not be in set A, so there is $1\cdot\binom{10}{2}$ (for the remaining elements).

 Case 4: Set A has 4 elements. Set B has 8 elements. 8 must be in set A, and 4 must not be in set A, so there is $1\cdot\binom{10}{3}$ (for the remaining element).

 Case 5: Set A has 5 elements. Set B has 7 elements. 7 must be in set A, and 5 must not be in set A, so there is $1\cdot\binom{10}{4}$ (for the remaining element).

 There is no Case 6 because $6$ can't go in set A and can't go in set B. Cases 7 to 11 are symmetric, so add up the 5 cases and multiply by $2$ to get $\boxed{772}$.

 -unhappyfarmer

 
                \begin{itemize}
                    \item Note by Voidling: After the first couple of cases, you may be able to find a pattern, where case n has $\binom{10}{n-1}$ possibilities, so you could just find $2\cdot\sum_{n=1}^{5}\binom{10}{n-1}$. This can be proved because for each case, if set A has $n$ elements, then out of these n elements, one of which must be $12-n$, there are 10 choices remaining, as none of the elements may equal $n$ either. Therefore, if A has n elements, then there are $\binom{10}{n-1}$ different possibilities for set A.
                \end{itemize}


                




	\item (\textit{AIME II 2010}) Let $ABCDEF$ be a regular hexagon. Let $G$, $H$, $I$, $J$, $K$, and $L$ be the midpoints of sides $AB$, $BC$, $CD$, $DE$, $EF$, and $FA$, respectively. The segments $\overline{AH}$, $\overline{BI}$, $\overline{CJ}$, $\overline{DK}$, $\overline{EL}$, and $\overline{FG}$ bound a smaller regular hexagon. Let the ratio of the area of the smaller hexagon to the area of $ABCDEF$ be expressed as a fraction $\frac {m}{n}$ where $m$ and $n$ are relatively prime positive integers. Find $m + n$.


 
 \textbf{Answer: }011



\textbf{Solution 1}

Without loss of generality, let $BC=2.$

 Note that $\angle BMH$ is the vertical angle to an angle of the regular hexagon, so it has a measure of $120^\circ$.

 Because $\triangle ABH$ and $\triangle BCI$ are rotational images of one another, we get that $\angle{MBH}=\angle{HAB}$ and hence $\triangle ABH \sim \triangle BMH \sim \triangle BCI$.

 Using a similar argument, $NI=MH$, and

 
 \[MN=BI-NI-BM=BI-(BM+MH).\]

 Applying the \href{https://artofproblemsolving.com/wiki/index.php?title=Law\_of\_cosines}{Law of cosines} on $\triangle BCI$, $BI=\sqrt{2^2+1^2-2(2)(1)(\cos(120^\circ))}=\sqrt{7}$

 
 \begin{align*}\frac{BC+CI}{BI}&=\frac{3}{\sqrt{7}}=\frac{BM+MH}{BH} \\ BM+MH&=\frac{3BH}{\sqrt{7}}=\frac{3}{\sqrt{7}} \\ MN&=BI-(BM+MH)=\sqrt{7}-\frac{3}{\sqrt{7}}=\frac{4}{\sqrt{7}} \\ \frac{\text{Area of smaller hexagon}}{\text{Area of bigger hexagon}}&=\left(\frac{MN}{BC}\right)^2=\left(\frac{2}{\sqrt{7}}\right)^2=\frac{4}{7}\end{align*}


 Thus, the answer is $4 + 7 = \boxed{011}.$

 





\textbf{Solution 2 (Coordinate Bash)}

We can use coordinates. Let $O$ be at $(0,0)$ with $A$ at $(1,0)$, 

 then $B$ is at $(\cos(60^\circ),\sin(60^\circ))=\left(\frac{1}{2},\frac{\sqrt{3}}{2}\right)$, 

 $C$ is at $(\cos(120^\circ),\sin(120^\circ))=\left(-\frac{1}{2},\frac{\sqrt{3}}{2}\right)$, 

 $D$ is at $(\cos(180^\circ),\sin(180^\circ))=(-1,0)$, 

 
 \begin{align*}&H=\frac{B+C}{2}=\left(0,\frac{\sqrt{3}}{2}\right) \\ &I=\frac{C+D}{2}=\left(-\frac{3}{4},\frac{\sqrt{3}}{4}\right)\end{align*}


 


 Line $AH$ has the slope of $-\frac{\sqrt{3}}{2}$and the equation of $y=-\frac{\sqrt{3}}{2}(x-1)$

 


 Line $BI$ has the slope of $\frac{\sqrt{3}}{5}$and the equation $y-\frac{3}{2}=\frac{\sqrt{3}}{5}\left(x-\frac{1}{2}\right)$

 
 

 Let's solve the system of equation to find $M$

 
 \begin{align*}-\frac{\sqrt{3}}{2}(x-1)-\frac{3}{2}&=\frac{\sqrt{3}}{5}\left(x-\frac{1}{2}\right) \\ -5\sqrt{3}x&=2\sqrt{3}x-\sqrt{3} \\ x&=\frac{1}{7} \\ y&=-\frac{\sqrt{3}}{2}(x-1)=\frac{3\sqrt{3}}{7}\end{align*}


 Finally,

 
 \begin{align*}&\sqrt{x^2+y^2}=OM=\frac{1}{7}\sqrt{1^2+(3\sqrt{3})^2}=\frac{1}{7}\sqrt{28}=\frac{2}{\sqrt{7}} \\ &\frac{\text{Area of smaller hexagon}}{\text{Area of bigger hexagon}}=\left(\frac{OM}{OA}\right)^2=\left(\frac{2}{\sqrt{7}}\right)^2=\frac{4}{7}\end{align*}


 Thus, the answer is $\boxed{011}$.

 
Diagram (by dragoon)\href{https://services.artofproblemsolving.com/download.php?id=YXR0YWNobWVudHMvZC82L2ZiNDljZmZiNjUzYWE4NGRmNmIwYTljMWQxZDU2ZDc1ZmNiMDQ3LmpwZWc=&rn=RDQ3ODA2RjUtMzlDNi00QzQ3LUE2OTYtMjlCQkE4NThDNkRBLmpwZWc=}{https://services.artofproblemsolving.com/download.php?id=YXR0YWNobWVudHMvZC82L2ZiNDljZmZiNjUzYWE4NGRmNmIwYTljMWQxZDU2ZDc1ZmNiMDQ3LmpwZWc=&rn=RDQ3ODA2RjUtMzlDNi00QzQ3LUE2OTYtMjlCQkE4NThDNkRBLmpwZWc=}

 





\textbf{Solution 3}

Use the diagram. Now notice that all of the "overlapping triangles" are congruent. You can use the AA similarity to see that the small triangles are similar to the large triangles. Now you can proceed as in Solution 1.

 




	\item (\textit{AIME II 2010}) Find the number of second-degree polynomials $f(x)$ with integer coefficients and integer zeros for which $f(0)=2010$.


 
 \textbf{Answer: }163



\textbf{Solution 1}

Let $f(x) = a(x-r)(x-s)$. Then $ars=2010=2\cdot3\cdot5\cdot67$. First consider the case where $r$ and $s$ (and thus $a$) are positive. There are $3^4 = 81$ ways to split up the prime factors between $a$, $r$, and $s$. However, $r$ and $s$ are indistinguishable. In one case, $(a,r,s) = (2010,1,1)$, we have $r=s$. The other $80$ cases are double counting, so there are $40$.

 We must now consider the various cases of signs. For the $40$ cases where $|r|\neq |s|$, there are a total of four possibilities, For the case $|r|=|s|=1$, there are only three possibilities, $(r,s) = (1,1); (1,-1); (-1,-1)$ as $(-1,1)$ is not distinguishable from the second of those three.

 You may ask: How can one of ${r, s}$ be positive and the other negative? $a$ will be negative as a result. That way, it's still $+2010$ that gets multiplied.

 Thus the grand total is $4\cdot40 + 3 = \boxed{163}$.

 Note: The only reason why we can be confident that $r = s$ is the only case where the polynomials are being overcounted is because of this: We have the four configurations listed below:

 $(a,r,s)\\ (a,-r,-s)\\ (-a,-r,s)\\ (-a,r,-s)$

 And notice, we start by counting all the positive solutions. So $r$ and $s$ must be strictly positive, no $0$ or negatives allowed. The negative transformations will count those numbers.

 So with these we can conclude that only the first and second together have a chance of being equal, and the third and fourth together. If we consider the first and second, the $x$ term would have coefficients that are always different, $-a(r + s)$ and $a(r + s)$ because of the negative $r$ and $s$. Since the $a$ is never equal, these can never create equal $x$ coefficients. We don't need to worry about this as $r$ and $s$ are positive and so that won't have any chance.

 However with the $(-a,-r,s)$ and $(-a,r,-s)$, we have the coefficients of the $x$ term as $a(s-r)$ and $a(r-s)$. In other words, they are equal if $s-r=r-s$ or $r=s$. Well if $r = 1$, then we have $s = 1$ and in the $(r,-s)$ case we have $(1,-1)$ and if we transform using $(s,-r)$, then we have $(-1, 1)$. So this is the only way that we could possibly overcount the equal cases, and so we need to make sure we don't count $(-1,1)$ and $(1,-1)$ twice as they will create equal sums. This is why we subtract $1$ from $41*4=164$.

 Each different transformation will give us different coordinates $(a,r,s)...$ it is just that some of them create equal coefficients for the $x$-term, and we see that they are equal only in this case by our exploration, so we subtract $1$ to account and get $163$.

 





\textbf{Solution 2}

We use \href{https://artofproblemsolving.com/wiki/index.php?title=Burnside\%27s\_Lemma}{Burnside's Lemma}. The set being acted upon is the set of integer triples $(a,r,s)$ such that $ars=2010$. Because $r$ and $s$ are indistinguishable, the permutation group consists of the identity and the permutation that switches $r$ and $s$. In cycle notation, the group consists of $(a)(r)(s)$ and $(a)(r \: s)$. There are $4 \cdot 3^4$ fixed points of the first permutation (after distributing the primes among $a$, $r$, $s$ and then considering their signs. We have 4 ways since we can keep them all positive, first 2 negative, first and third negative, or last two negative) and $2$ fixed points of the second permutation ($r=s=\pm 1$). By Burnside's Lemma, there are $\frac{1}{2} (4 \cdot 3^4+2)= \boxed{163}$ distinguishable triples $(a,r,s)$.Note: The permutation group is isomorphic to $\mathbb{Z}/2\mathbb{Z}$.

 





\textbf{Solution 3}

The equation can be written in the form of $k(x-a)(x-b)$, where $|k|$ is a divisor of $2010$. Factor $2010=2\cdot 3\cdot 5\cdot 67$. We divide into few cases to study. 

 Firstly, if $|k|=1$, the equation can be $-(x-a)(x+b)$ or $(x-a)(x-b)$ or $(x+a)(x+b)$, there are $2^4+2^4=32$ ways

 If $|k|\in \{2,3,5,67\}$. Take $|k|=2$ as an example, follow the procedure above, there are $2^3+2^3=16$ ways, and there are $\binom {4}{1}=4$ ways to choose $|k|$, so there are $16\cdot 4=64$ ways

 
If $|k|$ is the product of two factors of $2010$, there are $8$ ways for each. Thus there are $\binom {4}{2}\cdot 8=48$ ways

 If $|k|$ is the product of three factors of $2010$, there are $\binom {4}{3}\cdot 4=16$ ways

 In the end, $|k|=2010$, only $2010(x-1)(x-1); 2010(x+1)(x+1); -2010(x-1)(x+1)$ work, there are $3$ ways

 Thus, $32+64+48+16+3=\boxed{163}$

 





\textbf{Solution 4 (Vieta's)}

All polynomials are of the form $ax^2 + bx + 2010$.

 We have through Vieta's Formula $pq = \frac{2010}{a}$ so $a = \frac{2010}{pq}$. The condition $a \in \mathbb{Z}$ implies $pq$ is a divisor of 2010. So, $pq = 2^a3^b5^c67^d$. Either a factor of 2010 is in $p$, is in $q$, or is in neither $p$ nor $q$. The factor cannot be in both $p$ and $q$ as the resulting term is squared, and 2010 is square free. That is 81 viable cases. Now, either we have positive or negative roots, each contributing two cases for $(p,q)$ or $(q,p)$, giving us $4 \cdot 81 = 324$ plausible cases. If we were to divide by 2, as a polynomial with roots $(p,q)$ and $(q,p)$ are not at all distinct, then we have 162 solutions, so is the answer 162?

 No! We forgot to check for when $p = q$. You might say, "Pinotation, 2010 is squarefree!". However, if $p = q = 1, -1$, then such a solution $a = 2010$ still exists. These two orbits are unaffected under rotation, meaning we have to subtract it from the 324 ordered pairs, so we have $\frac{324-2}{2} + 2$ cases, or $161 + 2 = \boxed{163}$ plausible polynomials.

 This example of "unaffecting via rotation" is an example of a repeated (double) root, which is an example of a fixed orbit via Burnsides Lemma.

 




	\item (\textit{AIME II 2010}) Define a \textit{T-grid} to be a $3\times3$ matrix which satisfies the following two properties:


 
            \begin{enumerate}
                
\item Exactly five of the entries are $1$'s, and the remaining four entries are $0$'s.
\item Among the eight rows, columns, and long diagonals (the long diagonals are $\{a_{13},a_{22},a_{31}\}$ and $\{a_{11},a_{22},a_{33}\})$, no more than one of the eight has all three entries equal.
            \end{enumerate}
            Find the number of distinct \textit{T-grids}.


 
 \textbf{Answer: }068



\textbf{Solution 1}

The \textit{T-grid} can be considered as a tic-tac-toe board: five $1$'s (or X's) and four $0$'s (or O's). 

 There are only $\dbinom{9}{5} = 126$  ways to fill the board with five $1$'s and four $0$'s. Now we just need to subtract the number of bad grids. Bad grids are ones with more than one person winning, or where someone has won twice.

 Let three-in-a-row/column/diagonal be a "win" and let player $0$ be the one that fills in $0$ and player $1$ fills in $1$.

 \textbf{Case $1$:} Each player wins once. 

 If player X takes a diagonal, player Y cannot win. If either takes a row, all the columns are blocked, and visa versa. Therefore, they either both take a row or they both take a column.

 
            \begin{enumerate}
                <li>Both take a row: 
                \begin{itemize}
                    \item There are $3$ ways for player $1$ to pick a row, \item $2$ ways for player $0$ to pick a row and, \item $3$ ways for player $0$ to take a single box in the remaining row. 
                \end{itemize}


                Therefore, there are $3 \cdot 2 \cdot 3=18$ cases.

 </li>\item Both take a column: Using similar reasoning, there are $18$ cases.
            \end{enumerate}
            Case $1$: $36$ cases

 


 \textbf{Case $2$:} Player $1$ wins twice. (Player $0$ cannot win twice because he only has 4 moves.)

 
            \begin{enumerate}
                <li>Player $1$ picks a row and a column.
                \begin{itemize}
                    \item There are $3$ ways to pick the row\item and $3$ ways to pick the column. So there are $9$ cases.
                \end{itemize}


                </li><li>Player $1$ picks a row/column and a diagonal.
                \begin{itemize}
                    \item There are $6$ ways to pick the row/column \item and $2$ to pick the diagonal, 
                \end{itemize}


                so there are $12$ cases

 </li><li>2 diagonals
                \begin{itemize}
                    \item It is clear that there is only $1$ case. (Make a X).
                \end{itemize}


                </li>
            \end{enumerate}
            Case $2$ total: $22$

 


 Thus, the answer is $126-22-36=\boxed{68}$

 





\textbf{Solution 2}

We can use \href{https://artofproblemsolving.com/wiki/index.php?title=Generating\_functions}{generating functions} to compute the case that no row or column is completely filled with $1$'s. Given a row, let $a$, $b$, $c$ be the events that the first, second, third respective squares are $1$'s. Then the generating function representing the possible events that exclude a row of $1,1,1$ or $0,0,0$ from occuring is
 \[ab+bc+ca+a+b+c.\]Therefore, the generating function representing the possible grids where no row is filled with $0$'s and $1$'s is
 \[P(a,b,c)=((ab+bc+ca)+(a+b+c))^3.\]We expand this by the Binomial Theorem to find
 \[P(a,b,c)=(ab+bc+ca)^3+3(ab+bc+ca)^2(a+b+c)+3(ab+bc+ca)(a+b+c)^2+(a+b+c)^3.\]Recall that our grid has five $1$'s, hence we only want terms where the sum of the exponents is $5$. This is given by
 \[3(ab+bc+ca)^2(a+b+c).\]When we expand this, we find
 \[3(2abc(a+b+c)+a^2b^2+b^2c^2+c^2a^2)(a+b+c).\]We also want to make sure that each of $a$, $b$, $c$ appears at least once (so there is no column filled with $0$'s) and the power of each of $a$, $b$, $c$ is not greater than or equal to $3$ (so there is no column filled with $1$'s). The sum of the coefficients of the above polynomial is clearly $81$ (using $a,b,c=1$), and the sum of the coefficients of the terms $a^3bc$, $ab^3c$, and $abc^3$ is $6+6+6+3+3+3+3+3+3=36$, hence the sum of the coefficients of the desired terms is $81-36=45$. This counts the number of grids where no column or row is filled with	$0$'s or $1$'s. However, we could potentially have both diagonals filled with $1$'s, but this is the only exception to our $45$ possibilities, hence the number of $T$-grids with no row or column filled with the same digit is $44$.

 On the other hand, if a row (column) is filled with $0$'s, then by the Pigeonhole Principle, another row (column) must be filled with $1$'s. Hence this is impossible, so all other	possible $T$-grids have a row (column) filled with $1$'s. If the top row is filled with $1$'s, then we have two $1$'s left to place. Clearly they cannot go in the same row, because then the other row is filled with $0$'s. They also cannot appear in the same column. This leaves $3\cdot 2$ arrangements--3 choices for the location of the $1$ in the second row, and 2 choices for the location of the $1$ in the last row. However, two of these arrangements will fill a diagonal with $1$'s. Hence there are only $4$ $T$-grids where the top row is filled with $1$'s. The same argument applies if any other row or column is filled with $1$'s. Hence there are $4\cdot 6=24$ such $T$-grids.

 Thus the answer is $44+24=\boxed{68}$.

 





\textbf{Solution 3}

We proceed by complimentary counting.

 There are $\binom{9}{5} = 126$ ways to fill the board with five $1$'s and four $0$'s. We subtract the cases where at least two of the eight rows, columns, and long diagonals have all three entries equal. For the sake of simplicity, let a \textbf{line} denote a row/column/diagonal.

 Firstly, observe that there can be at most one line of $0$'s and at most two lines of $1$'s. Also, if there is one line of $0$'s, there can at most be one line of $1$'s. Therefore, there are only two "bad" cases we must consider:

 \textbf{Case $1$:} There is one line of $0$'s and one line of $1$'s.

 In this case, note that both lines must be horizontal or vertical. If both are horizontal, there are $3 \cdot 2 = 6$ ways to select the columns, and $\binom{3}{1} = 3$ ways to place the remaining $0$ in the remaining $3$ grids, yielding a total of $6 \cdot 3 = 18$ bad cases. Similarly, there are $18$ bad cases where both lines are vertical, so $18 + 18 = 36$ cases.

 \textbf{Case $2$:} There are two lines of $1$'s.

 In this case, note that all five $1$'s are used, and the two lines can be horizontal/vertical, horizontal/diagonal, vertical/diagonal, or diagonal/diagonal, yielding $3 \cdot 3 = 9$, $3 \cdot 2 = 6$, $3 \cdot 2 = 6$, and $1$ case, respectively. $9 + 6 + 6 + 1 = 22$ cases.

 


 Hence, the answer is $126 - 36 - 22 = \boxed{068}$.

 --OZ30

 




	\item (\textit{AIME II 2010}) Two noncongruent integer-sided isosceles triangles have the same perimeter and the same area. The ratio of the lengths of the bases of the two triangles is $8: 7$. Find the minimum possible value of their common perimeter.


 
 \textbf{Answer: }676



\textbf{Solution 1}

Let $s$ be the semiperimeter of the two triangles. Also, let the base of the longer triangle be $16x$ and the base of the shorter triangle be $14x$ for some arbitrary factor $x$. Then, the dimensions of the two triangles must be $s-8x,s-8x,16x$ and $s-7x,s-7x,14x$. By Heron's Formula, we have 

 <center>
 \[\sqrt{s(8x)(8x)(s-16x)}=\sqrt{s(7x)(7x)(s-14x)}\]
 \[8\sqrt{s-16x}=7\sqrt{s-14x}\]
 \[64s-1024x=49s-686x\]
 \[15s=338x\]

 </center>Since $15$ and $338$ are coprime, to minimize, we must have $s=338$ and $x=15$. However, we want the minimum perimeter. This means that we must multiply our minimum semiperimeter by $2$, which gives us a final answer of $\boxed{676}$.

 





\textbf{Solution 2}

Let the first triangle have sides $16n,a,a$, so the second has sides $14n,a+n,a+n$. The height of the first triangle is $\frac{7}{8}$ the height of the second triangle. Therefore, we have 
 \[a^2-64n^2=\frac{49}{64}((a+n)^2-49n^2).\] Multiplying this, we get 
 \[64a^2-4096n^2=49a^2+98an-2352n^2,\] which simplifies to 
 \[15a^2-98an-1744n^2=0.\] Solving this for $a$, we get $a=n\cdot\frac{218}{15}$, so $n=15$ and $a=218$ and the perimeter is $15\cdot16+218+218=\boxed{676}$.

 


 





\textbf{Solution 3}

Let the sides of the first triangle be ${a, a, 8k}$ and the sides of the second triangle be ${b, b, 7k}$

 Aim: We want to create some sort of relation between the perimeter and area since those elements are common in both triangles.

 First to find the perimeter of both triangles we do 

 $P = 2a + 8k$ where P is the perimeter of the first triangle$P = 2b + 8k$ where P is the perimeter of the second triangle

 Both equations use the same variable P since the perimeters of the triangles are equal as stated in the problem.

 Next to find the area of both triangles we multiply the height and base and divide by 2.

 To get the height of the first triangle we do pythag: $\sqrt{a^2 - (4k)^2}$Do the same for the second triangle and you get $\sqrt{b^2 - (7k/2)^2}$

 Multiple both by there bases and divide by 2 to get Area of 1st triangle = $(8k)^2*(\sqrt{4a^2-(8k)^2}$Area of 2nd triangle = $(7k)^2*(\sqrt{4b^2-(7k)^2}$

 Since both areas are equal we can create an equation:

 $(8k)^2*(\sqrt{4a^2-(8k)^2}$ =  $(7k)^2*(\sqrt{4b^2-(7k)^2}$

 Square both sides and divide by k\textasciicircum{}2:

 $64*(4a^2-64k^2) = 49(4c^2-49k^2)$

 Expand and combine like terms to get:

 
 \[256a^2 - 196c^2 = 1695k^2\]

 This factors into $(16a + 14c)(16a - 14c) = 1695k^2$

 recall that P = 2a + 8k and P = 2c + 7k, substituting those values gives

 
 \[(P - 15k)(15P - 113k) = 1695k^2\]

 Expand and combine like terms to get:

 
 \[15P^2 - 338Pk = 0\]

 We can further simplify this to:

 $P = \frac{338k}{15}$

 Proceed from this point like the other solutions!

 \textbf{Notes} The triangles in question have sides $\boxed{240,218,218}$ and $\boxed{210,233,233}$. ~eevee9406

 
We use $16x$ and $14x$ instead of $8x$ and $7x$ to ensure that the triangle has integral side lengths. Plugging $8x$ and $7x$ directly into Heron's gives $s=338$, but for this to be true, the second triangle would have side lengths of $\frac{223}{2}$, which is impossible. 

 




	\item (\textit{AIME II 2010}) The $52$ cards in a deck are numbered $1, 2, \cdots, 52$. Alex, Blair, Corey, and Dylan each pick a card from the deck randomly and without replacement. The two people with lower numbered cards form a team, and the two people with higher numbered cards form another team. Let $p(a)$ be the probability that Alex and Dylan are on the same team, given that Alex picks one of the cards $a$ and $a+9$, and Dylan picks the other of these two cards. The minimum value of $p(a)$ for which $p(a)\ge\frac{1}{2}$ can be written as $\frac{m}{n}$, where $m$ and $n$ are relatively prime positive integers. Find $m+n$.


 
 \textbf{Answer: }263



\textbf{Solution}

Once the two cards are drawn, there are $\dbinom{50}{2} = 1225$ ways for the other two people to draw. Alex and Dylan are the team with higher numbers if Blair and Corey both draw below $a$, which occurs in $\dbinom{a-1}{2}$ ways. Alex and Dylan are the team with lower numbers if Blair and Corey both draw above $a+9$, which occurs in $\dbinom{43-a}{2}$ ways. Thus, 
 \[p(a)=\frac{\dbinom{43-a}{2}+\dbinom{a-1}{2}}{1225}.\] Simplifying, we get $p(a)=\frac{(43-a)(42-a)+(a-1)(a-2)}{2\cdot1225}$, so we need $(43-a)(42-a)+(a-1)(a-2)\ge (1225)$. If $a=22+b$, then 
 \begin{align*}(43-a)(42-a)+(a-1)(a-2)&=(21-b)(20-b)+(21+b)(20+b)=2b^2+2(21)(20)\ge (1225) \\ b^2\ge \frac{385}{2} &= 192.5 >13^2 \end{align*}
 So $b> 13$ or $b< -13$, and $a=22+b<9$ or $a>35$, so $a=8$ or $a=36$. Thus, $p(8) = \frac{616}{1225} = \frac{88}{175}$, and the answer is $88+175 = \boxed{263}$.

 





\textbf{Solution 2}

Given that Alex and Dylan hold the cards $a$ and $a+9$, we need to calculate the probability that they end up on the same team. This happens in two scenarios:

 1. Both on the Lower Team: This occurs if the other two cards drawn are both greater than $a+9$.  2. Both on the Higher Team: This occurs if the other two cards drawn are both less than $a$.

 The total number of ways to choose the other two cards from the remaining 50 cards is $\binom{50}{2} = 1225$.

 The number of favorable outcomes is the sum of:  The number of ways to choose 2 cards greater than $a+9$: $\binom{43-a}{2}$  The number of ways to choose 2 cards less than $a$: $\binom{a-1}{2}$ 

 Thus, the probability $p(a)$ is:
 \[p(a) = \frac{\binom{a-1}{2} + \binom{43-a}{2}}{\binom{50}{2}} = \frac{\frac{(a-1)(a-2)}{2} + \frac{(43-a)(42-a)}{2}}{1225}.\]

 Finding the Minimum $a$ for $p(a) \geq \frac{1}{2}$

 Solving the inequality:
 \[\frac{\binom{a-1}{2} + \binom{43-a}{2}}{1225} \geq \frac{1}{2}\]leads to:
 \[2(a^2 - 44a + 291) \geq 0.\]

 This quadratic inequality is satisfied when $a \leq 8$ or $a \geq 36$.

 The minimum value of $p(a)$ that satisfies $p(a) \geq \frac{1}{2}$ occurs at $a = 8$ (or symmetrically at $a = 36$), giving:
 \[p(8) = \frac{616}{1225} = \frac{88}{175}\]where $88$ and $175$ are coprime.

 Final Answer
 \[88 + 175 = 263\]

 Answer: $263$

 




	\item (\textit{AIME II 2010}) Triangle $ABC$ with right angle at $C$, $\angle BAC < 45^\circ$ and $AB = 4$. Point $P$ on $\overline{AB}$ is chosen such that $\angle APC = 2\angle ACP$ and $CP = 1$. The ratio $\frac{AP}{BP}$ can be represented in the form $p + q\sqrt{r}$, where $p$, $q$, $r$ are positive integers and $r$ is not divisible by the square of any prime. Find $p+q+r$.


 
 \textbf{Answer: }007



\textbf{Solution 1}

Let $O$ be the \href{https://artofproblemsolving.com/wiki/index.php?title=Circumcenter}{circumcenter} of $ABC$ and let the intersection of $CP$ with the \href{https://artofproblemsolving.com/wiki/index.php?title=Circumcircle}{circumcircle} be $D$. It now follows that $\angle{DOA} = 2\angle ACP = \angle{APC} = \angle{DPB}$. Hence $ODP$ is isosceles and $OD = DP = 2$. 

 Denote $E$ the projection of $O$ onto $CD$. Now $CD = CP + DP = 3$. By the \href{https://artofproblemsolving.com/wiki/index.php?title=Pythagorean\_Theorem}{Pythagorean Theorem}, $OE = \sqrt {2^2 - \frac {3^2}{2^2}} = \sqrt {\frac {7}{4}}$. Now note that $EP = \frac {1}{2}$. By the Pythagorean Theorem, $OP = \sqrt {\frac {7}{4} + \frac {1^2}{2^2}} = \sqrt {2}$. Hence it now follows that,

 
 \[\frac {AP}{BP} = \frac {AO + OP}{BO - OP} = \frac {2 + \sqrt {2}}{2 - \sqrt {2}} = 3 + 2\sqrt {2}\]

 This gives that the answer is $\boxed{007}$.

 An alternate finish for this problem would be to use Power of a Point on $BA$ and $CD$. By Power of a Point Theorem, $CP\cdot PD=1\cdot 2=BP\cdot PA$. Since $BP+PA=4$, we can solve for $BP$ and $PA$, giving the same values and answers as above. 

 
                    \begin{figure}[H]
                        \centering
                        \includegraphics[width=0.5\textwidth]{Solution/AIME_II_2010_14_Solution_Figure_1.png}
                    \end{figure}
                    





\textbf{Solution 2}

Let $AC=b$, $BC=a$ by convention. Also, Let $AP=x$ and $BP=y$.  Finally, let $\angle ACP=\theta$ and $\angle APC=2\theta$. 

 We are then looking for $\frac{AP}{BP}=\frac{x}{y}$

 Now, by arc interceptions and angle chasing we find that $\triangle BPD \sim \triangle CPA$, and that therefore $BD=yb.$  Then, since $\angle ABD=\theta$ (it intercepts the same arc as $\angle ACD$) and $ADB$ is right, 

 $\cos\theta=\frac{DB}{AB}=\frac{by}{4}$.  

 
Using law of sines on $APC$, we additionally find that $\frac{b}{\sin 2\theta}=\frac{x}{\sin\theta}.$ Simplification by the double angle formula $\sin 2\theta=2\sin \theta\cos\theta$ yields 

 $\cos \theta=\frac{b}{2x}$.

 
We equate these expressions for $\cos\theta$ to find that $xy=2$. Since $x+y=AB=4$, we have enough information to solve for $x$ and $y$.  We obtain $x,y=2 \pm \sqrt{2}$ 

 Since we know $x>y$, we use $\frac{x}{y}=\frac{2+\sqrt{2}}{2-\sqrt{2}}=3+2\sqrt{2}$

 





\textbf{Solution 3}

Let $\angle{ACP}$ be equal to $x$. Then by Law of Sines, $PB = -\frac{\cos{x}}{\cos{3x}}$ and $AP = \frac{\sin{x}}{\sin{3x}}$. We then obtain $\cos{3x} = 4\cos^3{x} - 3\cos{x}$ and $\sin{3x} = 3\sin{x} - 4\sin^3{x}$. Solving, we determine that $\sin^2{x} = \frac{4 \pm \sqrt{2}}{8}$. Plugging this in gives that $\frac{AP}{PB} = \frac{\sqrt{2}+1}{\sqrt{2}-1} = 3 + 2\sqrt{2}$. The answer is $\boxed{007}$.

 
(You can derive that $\cos{3x} = 4\cos^3{x} - 3\cos{x},$ and similarly for $\sin{3x},$ by considering the expansion of $(\text{cis}(x))^3,$ equating real parts to $\cos{x}$ and imaginary parts to $\sin{x},$ then substituting with $1-\sin^2{x}$ to finish. ~happypi31415)

 





\textbf{Solution 4 (The quickest and most elegant)}

Let $\alpha=\angle{ACP}$, $\beta=\angle{ABC}$, and $x=BP$. By Law of Sines, 

 $\frac{1}{\sin(\beta)}=\frac{x}{\sin(90-\alpha)}\implies \sin(\beta)=\frac{\cos(\alpha)}{x}$ (1), and

 $\frac{4-x}{\sin(\alpha)}=\frac{4\sin(\beta)}{\sin(2\alpha)} \implies 4-x=\frac{2\sin(\beta)}{\cos(\alpha)}$. (2)

 Then, substituting (1) into (2), we get

 $4-x=\frac{2}{x} \implies x^2-4x+2=0 \implies x=2-\sqrt{2} \implies \frac{4-x}{x}=\frac{2+\sqrt{2}}{2-\sqrt{2}}=3+2\sqrt{2}$

 The answer is $\boxed{007}$.~Rowechen

 





\textbf{Solution 5}

Let $\angle{ACP}=x$. Then, $\angle{APC}=2x$ and $\angle{A}=180-3x$. Let the foot of the angle bisector of $\angle{APC}$ on side $AC$ be $D$. Then,

 $CD=DP$ and $\triangle{DAP}\sim{\triangle{APC}}$ due to the angles of these triangles.

 Let $CD=a$. By the Angle Bisector Theorem, $\frac{1}{a}=\frac{AP}{AD}$, so $AD=a\cdot{AP}$. Moreover, since $CD=DP=a$, by similar triangle ratios, $\frac{AP}{a+a\cdot{AP}}=a$. Therefore, $AP = \frac{a^2}{1-a^2}$. 

 Construct the perpendicular from $D$ to $AP$ and denote it as $F$. Denote the midpoint of $CP$ as $M$. Since $PD$ is an angle bisector, $PF$ is congruent to $PM$, so $PF=\frac{1}{2}$. 

 Also, $\triangle{DFA}\sim{\triangle{BCA}}$. Thus, $\frac{FA}{AC}=\frac{AD}{AB}\Longrightarrow\frac{\frac{a^2}{1-a^2}-\frac{1}{2}}{a+\frac{a^3}{1-a^2}}=\frac{\frac{a^3}{1-a^3}}{4}$. After some major cancellation, we have $7a^4-8a^2+2=0$, which is a quadratic in $a^2$. Thus, $a^2 = \frac{4\pm\sqrt{2}}{7}$. 

 Taking the negative root implies $AP<BP$, contradiction. Thus, we take the positive root to find that $AP=2+\sqrt{2}$. Thus, $BP=2-\sqrt{2}$, and our desired ratio is $\frac{2+\sqrt{2}}{2-\sqrt{2}}\implies{3+2\sqrt{2}}$.

 The answer is $\boxed{007}$.

 





\textbf{Solution 6}

Let $O$ be the circumcenter of $\triangle ABC$. Since $\triangle ABC$ is a right triangle, $O$ will be on $\overline{AB}$ and $\overline{AO} \cong \overline{OB} \cong \overline{OC} = 2$. Let $\overline{OP} = x$.

 Next, let's do some angle chasing. Label $\angle ACP = \theta^{\circ}$, and $\angle APC = 2\theta^{\circ}$. Thus, $\angle PAC = (180-3\theta)^{\circ}$, and by isosceles triangles,  $\angle ACO = (180-3\theta)^{\circ}$. Then, by angle subtraction, $\angle OCP = (\theta - (180-3\theta))^{\circ} = (4\theta - 180)^{\circ}$.

 Using the Law of Sines: 
 \[\frac{x}{\sin (4\theta-180)^{\circ}}=\frac{2}{\sin (2\theta)^{\circ}}\]Using trigonometric identies, $\sin (4\theta-180)^{\circ}=-\sin (4\theta)^{\circ}=-2\sin (2\theta)^{\circ}\cos (2\theta)^{\circ}$. Plugging this back into the Law of Sines formula gives us: 
 \[\frac{x}{-2\sin (2\theta)^{\circ}\cos (2\theta)^{\circ}}=\frac{2}{\sin (2\theta)^{\circ}}\]

 
 \[-4\sin (2\theta)^{\circ}\cos (2\theta)^{\circ}=x\sin (2\theta)^{\circ}\]
 \[-4\cos (2\theta)^{\circ}=x\]
 \[\cos(2\theta)^{\circ}=\frac{-x}4\]

 Next, using the Law of Cosines: 
 \[2^2=1^2+x^2-2\cdot 1\cdot x\cdot \cos (2\theta)^{\circ}\]Substituting $\cos(2\theta)^{\circ}=\frac{-x}4$ gives us:
 \[2^2=1^2+x^2-2\cdot 1\cdot x\cdot \frac{-x}4\]
 \[4=1+x^2+\frac{x^2}{2}\]

 Solving for x gives $x=\sqrt 2$

 Finally: $\frac{\overline{AP}}{\overline{BP}}=\frac{\overline{AO}+\overline{OP}}{\overline{BO}-\overline{OP}}=\frac{2+\sqrt 2}{2-\sqrt 2}=3+2\sqrt2$, which gives us an answer of $3+2+2=\boxed{007}$. ~adyj

 




	\item (\textit{AIME II 2010}) In triangle $ABC$, $AC=13$, $BC=14$, and $AB=15$. Points $M$ and $D$ lie on $AC$ with $AM=MC$ and $\angle ABD = \angle DBC$. Points $N$ and $E$ lie on $AB$ with $AN=NB$ and $\angle ACE = \angle ECB$. Let $P$ be the point, other than $A$, of intersection of the circumcircles of $\triangle AMN$ and $\triangle ADE$. Ray $AP$ meets $BC$ at $Q$. The ratio $\frac{BQ}{CQ}$ can be written in the form $\frac{m}{n}$, where $m$ and $n$ are relatively prime positive integers. Find $m-n$.


 
 \textbf{Answer: }218



\textbf{Solution 1 - Official Solution (MAA)}

The Angle Bisector Theorem implies that $E$ lies on $\overline{AN}$ and $D$ lies on $\overline{MC}$ because $AE/EB = AC/BC < 1$ and $AD/DC = AB/CB > 1$. The Angle Bisector Theorem furthermore implies
 \[NE=AN-AE=\frac 12\cdot AB - \frac{AC}{AC+BC}\cdot AB=\frac 5{18}\]and
 \[MD = CM - CD =  \frac 12\cdot AC - \frac{BC}{BC+BA}\cdot AC = \frac{13}{58}.\]Because $ANPM$ is cyclic, $\angle ENP = \angle ANP = \angle PMD$. Because $AEPD$ is cyclic, $\angle NEP = 180^\circ-\angle AEP = \angle MDP$. Because $\angle ENP =\angle PMD$ and $\angle NEP = \angle MDP$, triangles $NEP$ and $MDP$ are similar. Hence
 \[\frac{NE}{MD}=\frac{NP}{MP}.\]Applying the Law of Sines to $\triangle ANP$ and $\triangle AMP$ gives
 \[\frac{NE}{MD} = \frac{NP}{MP} = \frac{\sin  \angle NAP}{\sin  \angle PAM} = \frac{\sin  \angle BAQ}{\sin  \angle QAC}\]and thus
 \[\frac{\sin  \angle BAQ}{\sin  \angle QAC} = \frac{(\frac{5}{18})}{(\frac{13}{58})} = \frac{145}{117}.\]Thus
 \[\frac{BQ}{QC} = \frac{\textrm{Area}(ABQ)}{\textrm{Area}(ACQ)}  = \frac{\frac{1}{2}\cdot AB \cdot AQ \cdot \sin \angle BAQ}{\frac{1}{2}\cdot AC \cdot AQ \cdot \sin \angle QAC}= \frac{15}{13} \cdot \frac{145}{117} = \frac{725}{507},\]and $m - n = \boxed{218}$.

 





\textbf{Solution 2 (Linearity)}

Define the function $f:\mathbb{R}^{2}\rightarrow\mathbb{R}$ by 
 \[f(X)=\text{Pow}_{(AMN)}(X)-\text{Pow}_{(ADE)}(X)\] for points $X$ in the plane. Then $f$ is linear, so $\frac{BQ}{CQ}=\frac{f(B)-f(Q)}{f(Q)-f(C)}$. But $f(Q)=0$ since $Q$ lies on the radical axis of $(AMN)$, $(ADE)$ thus 
 \[\frac{BQ}{CQ}=-\frac{f(B)}{f(C)}=-\frac{BN\cdot BA-BE\cdot BA}{CM\cdot CA-CD\cdot CA}\] Let $AC=b$, $BC=a$ and $AB=c$. Note that $BN=\frac{c}{2}$ and $CM=\frac{b}{2}$ because they are midpoints, while $BE=\frac{ac}{a+b}$ and $CD=\frac{ab}{a+c}$ by Angle Bisector Theorem. Thus we can rewrite this expression as 
 \begin{align*}&-\frac{\frac{c^{2}}{2}-\frac{ac^{2}}{a+b}}{\frac{b^{2}}{2}-\frac{ab^{2}}{a+c}} \\ =&-\left(\frac{c^{2}}{b^{2}}\right)\left(\frac{\frac{1}{2}-\frac{a}{a+b}}{\frac{1}{2}-\frac{a}{a+c}}\right) \\ =&\left(\frac{c^{2}}{b^{2}}\right)\left(\frac{a+c}{a+b}\right) \\ =&\left(\frac{225}{169}\right)\left(\frac{29}{27}\right) \\ =&~\frac{725}{507}\end{align*}
 so $m-n=\boxed{218}$.

 





\textbf{Solution 3}

Let $Y = MN \cap AQ$. $\frac {BQ}{QC} = \frac {NY}{MY}$ since $\triangle AMN \sim \triangle ACB$. Since quadrilateral $AMPN$ is cyclic, $\triangle MYA \sim \triangle PYN$ and $\triangle MYP \sim \triangle AYN$, yielding $\frac {YM}{YA} = \frac {MP}{AN}$ and $\frac {YA}{YN} = \frac {AM}{PN}$. Multiplying these together yields *$\frac {YN}{YM} = \left(\frac {AN}{AM}\right) \left(\frac {PN}{PM}\right)$.

 $\frac {AN}{AM} = \frac {\frac {AB}{2}}{\frac {AC}{2}} = \frac {15}{13}$. 

 Now we claim that $\triangle PMD \sim \triangle PNE$. To prove this, we can use cyclic quadrilaterals.

 From $AMPN$, $\angle PNY \cong \angle PAM$ and $\angle ANM \cong \angle APM$. So, $m\angle PNA = m\angle PNY + m\angle ANM = m\angle PAM + m\angle APM = 180-m\angle PMA$ and $\angle PNA \cong \angle PMD$.

 From $ADPE$, $\angle PDE \cong \angle PAE$ and $\angle EDA \cong \angle EPA$. Thus, $m\angle MDP = m\angle PDE + m\angle EDA =  m\angle PAE + m\angle EPA = 180-m\angle PEA$ and $\angle PDM \cong \angle PEN$.

 Thus, from AA similarity, $\triangle PMD \sim \triangle PNE$.

 Therefore, $\frac {PN}{PM} = \frac {NE}{MD}$, which can easily be computed by the angle bisector theorem to be $\frac {145}{117}$. It follows that *$\frac {BQ}{CQ} = \frac {15}{13} \cdot \frac {145}{117} = \frac {725}{507}$, giving us an answer of $725 - 507 = \boxed{218}$.

 
                \begin{itemize}
                    \item These two ratios are the same thing and can also be derived from the Ratio Lemma.
                \end{itemize}


                Ratio Lemma :$\frac{BD}{DC} = \frac{AB}{AC} \cdot \frac{\sin \angle BAD}{\sin \angle CAD}$, for any cevian AD of a triangle ABC.For the sine ratios use Law of Sines on triangles APM and APN, 
 \[\frac{PM}{\sin \angle PAM}=\frac{AP}{\sin \angle AMP}=\frac{AP}{\sin \angle ANP}=\frac{PN}{\sin \angle PAN}\]. The information needed to use the Ratio Lemma can be found from the similar triangle section above.

 Source: \href{https://artofproblemsolving.comhttp://www.artofproblemsolving.com/Forum/viewtopic.php?p=1831745\#p1831745}{[1]} by Zhero

 





\textbf{Solution 4}

This problem can be solved with barycentric coordinates. Let triangle $ABC$ be the reference triangle with $A=(1,0,0)$, $B=(0,1,0)$, and $C=(0,0,1)$. Thus, $N=(1:1:0)$ and $M=(1:0:1)$. Using the Angle Bisector Theorem, we can deduce that $D=(14:0:15)$ and $E = (14:13:0)$. Plugging the coordinates for triangles $ANM$ and $AED$ into the circle formula, we deduce that the equation for triangle $ANM$ is $-a^2yz-b^2zx-c^2xy+(\frac{c^2}{2}y+\frac{b^2}{2}z)(x+y+z)=0$ and the equation for triangle $AED$ is $-a^2yz-b^2zx-c^2xy+(\frac{14c^2}{27}y+\frac{14b^2}{29}z)(x+y+z)=0$. Solving the system of equations, we get that $\frac{c^2y}{54}=\frac{b^2z}{58}$. This equation determines the radical axis of circles $ANM$ and $AED$, on which points $P$ and $Q$ lie. Thus, solving for $\frac{z}{y}$ gets the desired ratio of lengths, and $\frac{z}{y}=\frac{58c^2}{54b^2}$ and plugging in the lengths $b=13$ and $c=15$ gets $\frac{725}{507}$. From this we get the desired answer of $725-507=\boxed{218}$.-wertguk

 





\textbf{Solution 5 (Fast)}

Observe that $P$ is the center of spiral symmetry of segments $DM$ and $EN$. 

 Using the angle bisector theorem, we compute $DM = \frac{13}{58}$ and $EN = \frac{5}{18}$. Hence, it follows that the side-length ratio of triangles $DMP$ to $ENP$ is $\frac{117}{145}$ (note that the two triangles are similar by the spiral symmetry). This implies that the ratio of the height from $P$ to $AC$ to the height from $P$ to $AB$ is $\frac{117}{145}$, so we compute the area ratio $\frac{APC}{APB} = \frac{507}{725}$. 

 From the above, we see that the barycentric coordinates of $P$ are of the form $(x : 507 : 725)$. Hence, it follows that the point $Q$ has the coordinates $(0 : 507 : 725)$, so $\frac{BQ}{CQ} = \frac{725}{507}$ and our answer is $\boxed{218}$.

 \textbf{Extension} The work done in this problem leads to a nice extension of this problem:

 Given a $\triangle ABC$ and points $A_1$, $A_2$, $B_1$, $B_2$, $C_1$, $C_2$, such that $A_1$, $A_2$ $\in BC$, $B_1$, $B_2$ $\in AC$, and $C_1$, $C_2$ $\in AB$, then let $\omega_1$ be the circumcircle of $\triangle AB_1C_1$ and $\omega_2$ be the circumcircle of $\triangle AB_2C_2$. Let $A'$ be the intersection point of $\omega_1$ and $\omega_2$ distinct from $A$. Define $B'$ and $C'$ similarly. Then $AA'$, $BB'$, and $CC'$ concur. 

 This can be proven using Ceva's theorem and the work done in this problem, which effectively allows us to compute the ratio that line $AA'$ divides the opposite side $BC$ into and similarly for the other two sides.

 




\end{enumerate}